% 八、问题四
\section{问题四：波动电价下的购电策略}
问题 4 把电价变成随机过程，储能的角色从平滑负载扩展为价格套利\cite{hejazi}。附件 4 电价的全年逐段均值恰为附件 1 的常数电价曲线，即问题 2/3 是问题 4 的“平均电价”特例；波动电价一方面放大峰谷价差带来的套利空间，另一方面让“何时买、何时调”的判断依赖于\emph{还未实现}的价格。模型在前两问框架上只新增一个电价预测层，并按交易权限分为两支重算。

\subsection{电价预测层：日水平 $\times$ 日内形状}
电价预测文献\cite{weron}表明日内形状远比日间水平易于预测。电价与净负荷（负载减光伏）同受供需驱动，采用与负载预测同构的分解结构。对此前 $K_p=35$ 个完整日计算日价格水平 $\ell_i=T^{-1}\sum_tp_{i,t}$ 与平均归一化日内形状 $\chi_{d,t}$，两支分别建立水平模型：

\textbf{4-2 支（净负荷回归）}：价格水平由前一日水平与当日净负荷预测回归，
\begin{equation}
\hat\ell_{d\mid\tau}=a+\rho\,\ell_{d-1}+b\,\widehat N_{d\mid\tau}/10^5,\qquad
\hat p_{d,t\mid\tau}=\hat\ell_{d\mid\tau}\,\chi_{d,t},
\label{eq:q4price}
\end{equation}
其中 $\widehat N_{d\mid\tau}$ 为发布时刻 $\tau$ 的当日净负荷估计（已观测段真值 + 剩余段预测）。关键的因果细节：回归的\emph{历史}输入同样使用“历史当时的净负荷预测”重建，绝不用历史真实全天净负荷替代——否则拟合环境与使用环境不一致。\textbf{4-3 支}使用不含净负荷项的 AR(1) 水平模型。日内更新时按最近已完成段的预测误差做衰减修正；样本不足或回归退化按因果回退规则处理，预测价格下限 $10^{-4}$。

价格残差与负载、光伏残差按\textbf{同一历史发布时刻的完整路径}三通道配对（式\eqref{eq:q2scenario} 增加价格通道），保持“高负载日往往高电价”的跨通道相关。优化目标中：调整与普通购电费用（跨场景共享的第一阶段变量）按场景均价计费，各场景的紧急购电与跨午夜暂拟购电按该场景自身价格计费，实际结算始终使用附件 4 真值。

图~\ref{fig:fig_f10_price_forecast} 给出预测质量的两端：3 月 20 日（MAE 0.030 元/kWh）形状与水平俱佳；12 月 21 日（MAE 0.191 元/kWh）遭遇价格水平跳升，回归无法预见。全年逐日 MAE 的这种右偏分布正是问题 4 相对完美价格信息的费用差来源。

\begin{figure}[!htbp]\centering
\safeincludegraphics[width=0.92\textwidth]{fig_f10_price_forecast}
\caption{电价 0:00 预测与真值：左为 3 月 20 日（MAE 0.030 元/kWh，形状水平俱准），右为 12 月 21 日（MAE 0.191 元/kWh，水平跳升不可预见）——诚实展示预测的两端}
\label{fig:fig_f10_price_forecast}
\end{figure}

\subsection{两支的优化模型与权限差异}
\textbf{4-2（对应问题 2 权限）}：仅 0:00 冻结当天采购，不读取附件 3。其 0:00 优化模型与式\eqref{eq:q2saa} 同构，仅把常数电价换成价格场景：
\begin{equation}
\min_{G^0,\{\cdot_\omega\}}\ \sum_t\bar p_tG^0_t
+\frac1M\sum_\omega\Bigl(\sum_t5p_{\omega,t}E_{\omega,t}-v_\tau S_{\omega,144}\Bigr)
\quad\text{s.t. 式\eqref{eq:q2saa} 的全部约束},
\label{eq:q42}
\end{equation}
$\bar p_t$ 为场景均价，$v_\tau$ 为价格场景次日 0--5 时均价除以 $\eta$（随场景刷新，取代常数 0.481548）。虽无交易权限，但\textbf{价值更新不需要权限}：0/6/12/18 时仍重建价格预测、价格场景与 $\bar H$，执行层的保留水平式\eqref{eq:q2reserve} 使用当段\emph{真值}电价现算——价格信息随时间改善，执行随之变聪明，采购保持冻结。

\textbf{4-3（对应问题 3 权限）}：完整沿用式\eqref{eq:q3lp}，把各处电价换成对应的价格场景（调整费用按场景均价、场景内费用按各自价格），每个发布时刻同时刷新光伏预报、价格预测与剩余采购。

\begin{proposition}[价格同比缩放不变性]\label{prop:scale}
设全部价格输入同乘 $\lambda>0$：历史与真实电价 $p\mapsto\lambda p$（紧急电价保持 5 倍关系、调整费保持 $1.5$ 与 $0.5$ 倍关系），续存费率 $v\mapsto\lambda v$。则 (a) 各采购线性规划（问 1 确定性 LP、问 2 SAA、问 3 重优化）的最优解集不变、最优费用同乘 $\lambda$；(b) 价格预测层输出同乘 $\lambda$；(c) 未来费用函数 $\bar H_t\mapsto\lambda\bar H_t$，动态保留水平 $R_t$ 不变，DP 价值执行器的因果动作逐段不变；(d) 实际账单同乘 $\lambda$。即策略形状只由价格的相对结构决定，与价格水平无关。
\end{proposition}

\begin{proof}
(a) 三个线性规划的约束（能量平衡、库存递推、界）均不含价格；价格只出现在目标向量中，且各项费用（普通购电 $p$、调增 $1.5p$、调减 $-0.5p$、紧急 $5p$、续存 $-v$）对 $\lambda$ 一次齐次，故目标向量整体乘 $\lambda$，最优解集不变、最优值乘 $\lambda$。

(b) 价格预测为近 $K_p$ 个完整历史日的水平均值与归一化形状之积再乘水平因子（式\eqref{eq:q4price}）；水平、形状归一化与回归修正各环节均为历史价格的一次齐次函数，故输出乘 $\lambda$。

(c) $\bar H_t$ 由逐段下确界卷积生成：单段动作费用（固定采购段 $5p[r+\psi(x)]^+$、自由段 $p[L-V+\psi(x)]^+$）与终端项 $-\lambda ve$ 均对 $\lambda$ 一次齐次；下确界卷积与逐点平均都与正常数因子可交换，逆向归纳得 $\bar H_t\mapsto\lambda\bar H_t$。于是保留水平目标 $\lambda[\bar H_{t+1}(z)+5p_t\eta z]$ 的最小点集不变，$R_t$ 不变，执行动作逐段不变。

(d) 动作不变、结算价格乘 $\lambda$，账单各项一次齐次，总费用乘 $\lambda$。
\end{proof}


命题~\ref{prop:scale} 的数值验证（全局 0.5$\times$/2$\times$ 缩放）四项全部通过：LP 费用同比且原最优解仍属缩放后最优解集；SAA 计划购电量逐段不变；保留水平逐段不变；电价预测逐段同比。其含义是：问题 4 策略的成败取决于价格的\emph{形状与可预测性}，而非价格水平本身——把全年电价整体抬高一倍不改变任何决策，改变决策的是“今天下午比明天凌晨贵多少、这个差能否被预见”。

\subsection{结果与解读}
主方法全年实际费用：\textbf{4-2 为 1422.99 万元，4-3 为 1369.30 万元}（24h 基线分别为 1427.06 与 1371.70 万元）。三个定位参照：附件 4 电价下的完美信息离线下界为 1278.09 万元；同一算法、完美价格信息（0:00 已知当天与次日真值电价、无电价残差场景，与 24h 基线同口径）的参照为 4-2 支 1418.18 万元、4-3 支 1362.34 万元。两个观察值得展开：

\textbf{其一，价格预测误差的代价意外地小}。同口径比较，4-2 基线与其完美价格参照差 8.88 万元（0.62\%），4-3 基线差 9.36 万元（0.68\%）。原因在结构上：采购与执行依赖的是价格的\emph{相对形状}（何时便宜何时贵），而日内形状 $\chi_{d,t}$ 高度稳定、易于预测；难以预测的是日间水平跳动，而命题~\ref{prop:scale} 恰说明水平误差对决策形状几乎无影响，只影响对少数极端日的调整幅度。

\textbf{其二，4-3 相对 4-2 的 53.69 万元差距与问题 3 相对问题 2 的 54.01 万元几乎一致}。两个差距都主要来自“正式预报 + 日内调整权限”这一组合，说明该组合的价值对电价环境稳健；波动电价没有放大也没有侵蚀调整权限的价值。逐月看（图~\ref{fig:fig_f11_q4_monthly}），4-3 的优势集中在光伏波动大的月份，2 月甚至出现小幅负值——当月负载水平低、调整机会少，而 4-3 的 AR(1) 价格模型在当月水平跳变日恰好方向出错；诚实呈现该负值月份正说明差距的来源是信息而非口径。

\begin{figure}[!htbp]\centering
\safeincludegraphics[width=0.92\textwidth]{fig_f11_q4_monthly}
\caption{问题四两支相对各自常数电价对应问（4-2 对问 2、4-3 对问 3）的逐月费用差：波动电价整体抬升费用约 5\%，抬升幅度冬季最大；4-3 对 4-2 的月度优势与问 3 对问 2 的模式一致}
\label{fig:fig_f11_q4_monthly}
\end{figure}

四个分支处于不同的信息与交易环境，总费用只在各自环境内评价；4-2 与 4-3 之间同时改变了预测规则、可得光伏信息与采购权限，其差额是组合价值，本文不对其做单因素归因。完整策略已存入 result4-2.xlsx 与 result4-3.xlsx。
